\section{The Aldous--Lyons track: $\mathrm{TMIP}^* = \RE$}
\label{ch:tailored}

This chapter is the plan and the progress of a second target: the main theorem of Bowen, Chapman
and Vidick's \emph{The Aldous--Lyons Conjecture II: Undecidability}~\cite{BCV25}, a computable
reduction from the halting problem to \emph{tailored} non-local games that sends halting
machines to games with a perfect Z-aligned permutation strategy commuting along edges, and
non-halting machines to games of value at most $\frac12$. With the companion
paper~\cite{BCLV24} it refutes the Aldous--Lyons conjecture. The estimate, the reuse map and
the phases are in \texttt{planning/aldous-lyons-track.md}. There is no adversarial-verification
ledger for these papers: the paper, at its arXiv version, is the authority on the mathematics,
and line numbers II:$n$ below refer to its LaTeX source.

\emph{Status.} Phase~0: the definitions and the statement are formalized
(\texttt{MIPRE/TailoredGameValue.lean} and \texttt{MIPRE/Tailored/}), and one lemma is proved,
that the definitions are not vacuous (Lemma~\ref{lem:tailored-trivial-zpc}). The other results
below carry no proof marks; they are the targets of Phases~1--5.

The formalization departs from the paper's text in the places Remark~\ref{rem:tailored-route}
lists, all of them choices this project made for $\MIPstar = \RE$ already: soundness is proved
in value form, against all finite-dimensional strategies, with direct parallel repetition, so
that neither the anchored repetition theorem of~\cite{BVY17} nor the almost-synchronous
rounding of~\cite{Vid22} is needed.

\subsection{Tailored games and permutation strategies}

\begin{definition}[Tailored game; II:1243]
\label{def:tailored-game}
\uses{def:game,def:q-value,def:normal-verifier}
\lean{MIPRE.Tailored.TailoredGame, MIPRE.Tailored.TailoredGame.len,
  MIPRE.Tailored.TailoredGame.maxLen, MIPRE.Tailored.TailoredGame.Accepts,
  MIPRE.Tailored.TailoredGame.toGame, MIPRE.Tailored.TailoredGame.valStar,
  MIPRE.Tailored.TailoredGame.doubled, MIPRE.Tailored.Satisfies,
  MIPRE.Tailored.rejectConstraint}
\leanok
A \emph{tailored game} on a finite question set $X$ is a distribution $\mu$ on $X \times X$
together with, at each question $x$, a set $S_x = S_x^R \cup S_x^L$, a disjoint union, of $\ell^R(x)$
\emph{readable} and $\ell^L(x)$ \emph{linear} formal variables, and, for each pair of questions
$(x, y)$ and each assignment $\gamma^R$ of the readable variables at $x$ and $y$, a finite list
$L_{xy}(\gamma^R)$ of vectors over the disjoint union $S_x \cup S_y \cup \{J\}$, the \emph{controlled linear
constraints}. The answer at $x$ is an assignment $a \in \F_2^{S_x}$, readable bits first. The
\emph{canonical decider} (II:1757--1787) accepts $(a, b)$ at $(x, y)$ when every
$c \in L_{xy}(a^R, b^R)$ satisfies $\langle c, (a, b, 1)\rangle = 0$ over $\F_2$, the
coordinate $J$ being set to $1$; the empty list accepts every answer pair, and the list
$\{J\}$ rejects every one. The tailored game defines a game (Definition~\ref{def:game}) with
answers the bit strings of bounded length, those of the wrong length losing, and its quantum
value $\valstar$. Its \emph{doubled} game asks $(\mathsf{false}, x)$ and $(\mathsf{true}, y)$
with the weight of $(x, y)$ and no weight elsewhere; it is again a tailored game.
\end{definition}

\begin{definition}[Z-aligned permutation strategies commuting along edges; II:1043--1057,
II:1279--1283]
\label{def:zpc-strategy}
\uses{def:tailored-game}
\lean{MIPRE.Tailored.signedPermMatrix, MIPRE.Tailored.IsSignedPerm,
  MIPRE.Tailored.PermStrategy, MIPRE.Tailored.PermStrategy.proj,
  MIPRE.Tailored.PermStrategy.value, MIPRE.Tailored.TailoredGame.HasPerfectZPC}
\leanok
A \emph{signed permutation matrix} on $\C^m$ is a matrix $e_j \mapsto (-1)^{s(j)} e_{\sigma(j)}$
for a permutation $\sigma$ of $\{1, \dots, m\}$ and signs $s$: the action of a signed
permutation of $\{\pm\} \times \{1, \dots, m\}$ on the anti-symmetric functions, in their
standard basis. A \emph{permutation strategy} for a tailored game assigns to each variable $X$
at each question $x$ an observable $U(x, X)$ on $\C^m$ that is a signed permutation matrix and
an involution, the observables at a question commuting with each other; its measurement at $x$
is the Fourier transform
\[
  P^x_a \;=\; \prod_{X \in S_x} \tfrac12\bigl(\Id + (-1)^{a_X}\, U(x, X)\bigr),
  \qquad a \in \F_2^{S_x},
\]
and its value is $\sum_{x, y} \mu(x, y) \sum_{a, b} \mathrm{tr}(P^x_a P^y_b)/m$ over the
accepted $(a, b)$. It is \emph{Z-aligned} when the observables of readable variables are
diagonal, and \emph{commutes along edges} when $U(x, X)$ and $U(y, Y)$ commute whenever
$\mu(x, y) > 0$. A ZPC strategy is a Z-aligned permutation strategy commuting along edges, and
it is \emph{perfect} when its value is $1$.
\end{definition}

\begin{remark}[Two conventions]
\label{rem:tailored-conventions}
The paper's loops $xx$ read a single answer, over $S_x$; here, as for the paper's tailored
normal form verifiers, whose canonical decider always receives two answers, a loop reads the
pair $(a, a)$ against constraints over two copies of $S_x$ --- a one-copy constraint $(c, c_J)$
is the two-copy constraint $(c, 0, c_J)$, and a two-copy constraint $(c_1, c_2, c_J)$ acts as
$(c_1 + c_2, c_J)$. And the completeness clauses of the pipeline are read on the doubled game,
as Definition~\ref{def:pcc} is read for the existing pipeline: the doubled game is bipartite,
so no hypothesis on the verifier is needed to make it synchronous. A ZPC strategy for a game
gives one for its doubled game, the same observables at both tags.
\end{remark}

\begin{lemma}[The trivial strategy; II:1148]
\label{lem:tailored-trivial-zpc}
\uses{def:zpc-strategy}
\lean{MIPRE.Tailored.PermStrategy.trivial, MIPRE.Tailored.PermStrategy.trivial_proj,
  MIPRE.Tailored.hasPerfectZPC_of_accepts_zero}
\leanok
A tailored game that accepts the all-zero answers at every question pair of positive weight
has a perfect ZPC strategy: in dimension $1$, every observable the identity.
\end{lemma}

\begin{proof}
\leanok
Every observable is the identity, which is a signed permutation matrix, diagonal and an
involution, and commutes with everything. Each factor $\frac12(\Id + (-1)^{a_X}\Id)$ is $\Id$
at $a_X = 0$ and $0$ at $a_X = 1$, so $P^x_a$ is $1$ at $a = 0$ and $0$ elsewhere, and the
value is $\sum_{x, y} \mu(x, y) = 1$. This is the case of the halting protocol in which the
machine has halted (II:1978), and it shows that the definitions are not vacuous.
\end{proof}

\begin{lemma}[ZPC strategies are PCC strategies]
\label{lem:zpc-pcc}
\uses{def:zpc-strategy,def:pcc,def:sync-value}
The measurements $P^x$ of a permutation strategy form a synchronous strategy
(Definition~\ref{def:sync-strategy}) of the same value, commuting on the support of $\mu$
when the permutation strategy commutes along edges. In particular a tailored game with a
perfect ZPC strategy has synchronous value $1$, and quantum value $1$.
\end{lemma}

\begin{proof}
The observables at $x$ are commuting involutions, so the factors
$\frac12(\Id \pm U(x, X))$ are commuting projections, $P^x_a$ is a projection, and the
$P^x_a$ sum to $\Id$ by expanding the product. Commutation of the observables gives
commutation of the measurements, and the two values are the same sum. Phase~1.
\end{proof}

\subsection{Tailored normal form verifiers}

\begin{definition}[Tailored normal form verifier; II:1758]
\label{def:tailored-verifier}
\uses{def:tailored-game,def:zpc-strategy,def:sampler,def:decider}
\lean{MIPRE.Tailored.TailoredVerifier, MIPRE.Tailored.LenIs, MIPRE.Tailored.LenBound,
  MIPRE.Tailored.LenTotal, MIPRE.Tailored.LpIs, MIPRE.Tailored.TailoredVerifier.progs,
  MIPRE.Tailored.TailoredVerifier.size, MIPRE.Tailored.TailoredVerifier.Questions,
  MIPRE.Tailored.TailoredVerifier.lenOf, MIPRE.Tailored.TailoredVerifier.LenDefined,
  MIPRE.Tailored.TailoredVerifier.consOf, MIPRE.Tailored.TailoredVerifier.tgame,
  MIPRE.Tailored.TailoredVerifier.valStar, MIPRE.Tailored.TailoredVerifier.HasPerfectZPC}
\leanok
A \emph{tailored normal form verifier} $\verifier = (\cal{S}, \cal{L}, \cal{P})$
is a conditionally linear sampler $\cal{S}$ (Definition~\ref{def:sampler}), an
\emph{answer-length calculator} $\cal{L}$, which on $(n, x, \kappa)$ outputs the number of
readable ($\kappa = R$) or linear ($\kappa = L$) variables at the question $x$, and a
\emph{linear-constraints processor} $\cal{P}$, which on $(n, x, y, a^R, b^R)$ outputs the
list $L_{xy}(a^R, b^R)$; the paper's fourth machine, the canonical decider, is the same for
every verifier and is left implicit. The \emph{$n$-th game} $\verifier_n$ is the tailored game
on the sampler's questions that the programs compute, a question or pair at which a program
does not halt rejecting every answer. $\verifier$ \emph{has a perfect ZPC strategy at $n$} when
the doubled game of $\verifier_n$ has one (Remark~\ref{rem:tailored-conventions}).
\end{definition}

\begin{definition}[$\lambda$-bounded tailored verifiers, and budgets; II:1800]
\label{def:tailored-bounded}
\uses{def:tailored-verifier}
\lean{MIPRE.Tailored.TailoredVerifier.IsBounded, MIPRE.Tailored.TailoredVerifier.Within}
\leanok
$\verifier$ is \emph{$\lambda$-bounded} when, for every $n \geq 2$, the sampler's dimension and
the running times of the three programs are at most $n^\lambda$ --- the running times read as
in Definition~\ref{def:lambda-bounded}, at degree $\lambda$ in the size of the input --- the
lengths $\cal{L}$ outputs are at most $n^\lambda$, and $\abs{\verifier} \leq \lambda$. The
paper's answer-length calculator outputs in unary, so that its running time bounds the
lengths; here the lengths are numbers and the bound is a clause of its own. A verifier is
\emph{within a budget} at index $n$ when the same quantities are bounded by the budget's.
\end{definition}

\subsection{The target}

\begin{definition}[Tailored game descriptions, standalone; II:1243, II:1279]
\label{def:tailored-game-data}
\uses{def:game-description}
\lean{TailoredGameValue.TailoredGameData, TailoredGameValue.TailoredGameData.equivTuple,
  TailoredGameValue.TailoredGameData.lenRAt, TailoredGameValue.TailoredGameData.lenLAt,
  TailoredGameValue.TailoredGameData.lenAt, TailoredGameValue.TailoredGameData.ansLen,
  TailoredGameValue.TailoredGameData.bit, TailoredGameValue.TailoredGameData.WellFormatted,
  TailoredGameValue.TailoredGameData.readable, TailoredGameValue.TailoredGameData.full,
  TailoredGameValue.TailoredGameData.Satisfies, TailoredGameValue.TailoredGameData.Accepts,
  TailoredGameValue.TailoredGameData.weights, TailoredGameValue.TailoredGameData.toGame,
  TailoredGameValue.signedPermMatrix, TailoredGameValue.PermStrategy,
  TailoredGameValue.PermStrategy.proj, TailoredGameValue.PermStrategy.value,
  TailoredGameValue.TailoredGameData.HasPerfectZPC}
\leanok
A \emph{tailored game description} lists the number of vertices, the readable and linear
lengths at each vertex, the question weights, and the controlled linear constraints as
entries $(x, y, \gamma^R, c)$, read as $c \in L_{xy}(\gamma^R)$. It denotes a synchronous game
in the sense of Definition~\ref{def:game-description}: questions are vertices, answers bit
vectors of the maximal length that vanish beyond the vertex's length, the weights are
normalized as there, and acceptance is the canonical decider's, unequal answers at a loop
losing. ZPC strategies are Definition~\ref{def:zpc-strategy}'s, the generators beyond a
vertex's length acting as the identity. These definitions are written on top of Mathlib alone,
in \texttt{MIPRE/TailoredGameValue.lean}, beside the standalone statement of
Theorem~\ref{thm:main}, whose game value they reuse.
\end{definition}

\begin{theorem}[$\mathrm{TMIP}^* = \RE$; II:1505, Theorem~7.4 of~\cite{BCLV24}]
\label{thm:tmip-re}
\uses{def:tailored-game-data}
\lean{TailoredGameValue.TailoredHaltingReduction}
\leanok
There is a computable map $g$ from Turing machines to tailored game descriptions such that for
every Turing machine $\machine$:
\begin{enumerate}
\item (\textbf{Completeness}) if $\machine$ halts on the empty input, then $g(\machine)$ has a
  perfect ZPC strategy;
\item (\textbf{Soundness}) if $\machine$ does not halt on the empty input, then
  $\synval(g(\machine)) \leq \frac12$.
\end{enumerate}
Turing machines, halting and computability are as in Theorem~\ref{thm:main}, and $\synval$ is
the synchronous value of Definition~\ref{def:sync-value}, in its standalone form
(Remark~\ref{rem:standalone-restatement}).
\end{theorem}

\begin{proof}
\uses{thm:tailored-halting,thm:tailored-compression-target}
Theorem~\ref{thm:tailored-halting} at the compression of
Theorem~\ref{thm:tailored-compression-target}, then the tabulation of the halting verifier's
game at its fixed index as a tailored game description, as Theorem~\ref{thm:main} tabulates a
normal form verifier. Phases~1 and~5.
\end{proof}

\begin{remark}[Where the formalization departs from the paper]
\label{rem:tailored-route}
\begin{enumerate}
\item \emph{The soundness bound.} The paper's clause~(3) reads $\valstar(\game_\machine) <
  \frac12$; its proof (II:2046--2065) shows that $\Ent(\game_\machine, \frac12) = \infty$, that
  every finite-dimensional strategy has value $< \frac12$, which bounds the supremum by
  $\frac12$ only. Paper~I uses the theorem with $\leq \frac12$ (\cite{BCLV24}, Theorem~7.4).
\item \emph{Value form.} The paper states compression and parallel repetition in entanglement
  form, $\Ent(\verifier'_n, \frac12) \geq \max\{\Ent(\verifier_{2^n}, \frac12),
  2^{2^{\lambda n} - 1}\}$ (II:1865), proved through the anchored repetition theorem
  of~\cite{BVY17} read as dimension-preserving (II:11425) and the almost-synchronous rounding
  of~\cite{Vid22} (II:3227, a proof idea; its quotation of~\cite{Vid22} at II:11350 is commented
  out of the source). The paper's Remark~II:1876 says that what its proofs establish is the
  value form, which is what is formalized here, with direct repetition
  (Theorem~\ref{thm:direct-repetition-q}) and the paper route of Theorem~\ref{thm:halting}, whose
  search branch replaces the base case $2^{2^{\lambda n} - 1}$. No entanglement bound is
  tracked.
\item \emph{All strategies.} Soundness is proved in the quantum value $\valstar$ of tensor-product
  strategies, which bounds the paper's synchronous value from above
  (Lemma~\ref{lem:sync-le-valstar}).
\item \emph{Question reduction.} The paper's question reduction is JNVWY's on its own Pauli basis
  game over $\F_2$ (II:3742), sound by a semi-stability theorem cited from de la Salle (II:3495);
  the paper notes that its combinatorial game is JNVWY's and that JNVWY's soundness covers it
  (II:5228--5233). The plan's Route~A presents this project's introspection verifier
  (Theorem~\ref{thm:introspection}) as a tailored verifier instead, which Phase~0 checked is
  possible (Remark~\ref{rem:tailored-route-a}).
\end{enumerate}
\end{remark}

\begin{remark}[Route A: the two checks]
\label{rem:tailored-route-a}
Two properties of this project's introspection verifier make the tailored presentation
possible, and both hold (\texttt{planning/aldous-lyons-track.md}, section~4.2, with the evidence).
Every check of the introspective predicate is \emph{controlled-linear}: a predicate on the
readable fields alone (the introspected question, the sampled seed, the prefixes read and
hidden, and the $Z$-basis Pauli answers), or $\F_2$-affine equations on the linear fields (the
duals, the hidden tails, the $X$-basis answers, the Pair and Magic Square answers) with
coefficients computed from readable data and the questions; the input's answers must first be
padded to constant lengths, as in the paper (II:6219). And the honest strategy is a
signed-permutation strategy with diagonal readable observables when the input's is: its
$X$-type observables are translations, its $Z$-type observables $\pm 1$ characters, its
nonlinear post-processing acts on $Z$-basis outcomes only, and the input enters as controlled
direct sums.
\end{remark}

\subsection{Compression and its three stages}

\begin{definition}[Tailored gap compression, value form; II:1852]
\label{def:tailored-compression}
\uses{def:tailored-bounded}
\lean{MIPRE.Tailored.TailoredGapCompression}
\leanok
A \emph{tailored gap compression} at level $\ell$ is a polynomial-time procedure that maps a
tailored verifier and $\lambda$ to a tailored verifier whose sampler and answer-length
calculator depend on $\lambda$ only, the calculator halting on every question of the sampler,
with running times, dimension and lengths at most a polynomial in $n + \lambda$, such that for
every $\lambda$-bounded input and every $n \geq C_0$: a perfect ZPC strategy for
$\verifier_{2^n}$ gives one for $\verifier'_n$; and $\valstar(\verifier_{2^n}) \leq \frac12$
implies $\valstar(\verifier'_n) \leq \frac12$.
\end{definition}

\begin{theorem}[Halting from tailored compression; II:1887--2066]
\label{thm:tailored-halting}
\uses{def:tailored-compression,def:tailored-game-data,lem:zpc-pcc,thm:halting}
A tailored gap compression gives the conclusion of Theorem~\ref{thm:tmip-re}.
\end{theorem}

\begin{proof}
The paper's halting protocol is the paper route of Theorem~\ref{thm:halting}, with the
linear-constraints processor in the role of the decider: the processor of $\verifier^{\machine,
\lambda}$ is a fixed point that outputs no constraint once $\machine$ has halted within the
step budget, rejects (the list $\{J\}$) when the search for $\valstar > \frac12$ at the fixed
level succeeds, and otherwise runs the compressed processor of its own description. When
$\machine$ halts, an always-accepting game has the deterministic ZPC strategy ($m = 1$, every
observable the identity) --- this is where the calculator's totality is spent --- and
completeness descends through the levels; when it does not, the downward induction of the
search branch gives $\valstar \leq \frac12$, and Lemma~\ref{lem:sync-le-valstar} the synchronous
value. Phase~1.
\end{proof}

\begin{definition}[The three stage contracts]
\label{def:tailored-stages}
\uses{def:tailored-bounded,def:introspection-contract,def:answer-reduction-contract,
  thm:parallel-repetition}
\lean{MIPRE.Tailored.TailoredIntrospection, MIPRE.Tailored.TailoredAnswerReduction,
  MIPRE.Tailored.TailoredRepetition, MIPRE.Tailored.TailoredRepetition.arg}
\leanok
Question reduction (II:5712), answer reduction (II:6883) and parallel repetition (II:11103) for
tailored verifiers, each in value form, with the parameters, budgets and soundness losses of
the corresponding contracts of the existing pipeline
(Definitions~\ref{def:introspection-contract} and~\ref{def:answer-reduction-contract},
Theorem~\ref{thm:parallel-repetition}): tailored verifiers in and out, the output's
answer-length calculator part of the data, completeness read in perfect ZPC strategies, and
repetition's soundness loss depending on a bound on the input's lengths in place of a parse
length.
\end{definition}

\begin{theorem}[Question reduction for tailored verifiers; II:5712]
\label{thm:tailored-qr}
\uses{def:tailored-stages,thm:introspection}
The contract \texttt{TailoredIntrospection} is inhabited, at input level~$7$.
\end{theorem}

\begin{proof}
Route~A (Remark~\ref{rem:tailored-route-a}): pad the input to constant lengths (II:6219),
present the introspection verifier of Theorem~\ref{thm:introspection} as a tailored verifier
whose constraints accept exactly what the introspective decider accepts, transport its
soundness through the identification of the games, and prove ZPC completeness of its honest
strategy. Phase~3.
\end{proof}

\begin{theorem}[Answer reduction for tailored verifiers; II:6883]
\label{thm:tailored-ar}
\uses{def:tailored-stages}
The contract \texttt{TailoredAnswerReduction} is inhabited, at input level~$5$.
\end{theorem}

\begin{proof}
The paper's tailored PCP (II:7736--11074): purification, oracularization, triangulation and
decoupling of the linear system, the output indicator and its decoupled succinct description,
the degree-$9$ PCP whose linear part is affine in the linear answers, and the seeded low-degree
test, which is this project's (\texttt{MIPRE.LIDT.CL.clGame}, with its soundness). Phase~4.
\end{proof}

\begin{theorem}[Parallel repetition of tailored verifiers; II:11103, direct form]
\label{thm:tailored-rep}
\uses{def:tailored-stages,thm:direct-repetition-q}
The contract \texttt{TailoredRepetition} is inhabited, at every level.
\end{theorem}

\begin{proof}
The $k$-fold tailored product (II:11377), its constraints padded with zeros to the whole answer;
ZPC completeness through Kronecker products of signed permutations (II:3043); soundness from
Theorem~\ref{thm:direct-repetition-q}. Phase~2.
\end{proof}

\begin{theorem}[Tailored compression; II:5643, value form]
\label{thm:tailored-compression-target}
\uses{def:tailored-compression,thm:tailored-qr,thm:tailored-ar,thm:tailored-rep}
There is a tailored gap compression at level~$7$.
\end{theorem}

\begin{proof}
Question reduction, then answer reduction, then repetition, the parameters chosen and the
budgets chained as in Theorem~\ref{thm:compression}, with a third program in every accounting
line. Phase~5.
\end{proof}

\begin{remark}[The Aldous--Lyons conjecture]
\label{rem:tailored-al}
Paper~I~\cite{BCLV24} associates with a tailored game a subgroup test whose sofic value is $1$
when the game has a perfect ZPC strategy and at most $1 - (C\Lambda^4 2^{6\Lambda})^{-1}$ when the
synchronous value of the game is at most $\frac12$, $\Lambda$ being the largest answer length
(Theorem~7.3 there, a finitary argument), and shows that under the Aldous--Lyons conjecture the
sofic value of a subgroup test is computable to any precision (Main Theorem~I there, a weak-$*$
compactness argument). With Theorem~\ref{thm:tmip-re} the conjecture is false. This is Phase~6
of the plan, after the main theorem.
\end{remark}
